\section{Current-device integration and research discovery}
\label{chapter:07}

% block-id: auto.07.conda
\hypertarget{auto.07.conda}{}
The integration on 28 October 2026 separates project, ANSYS and retrieval dependencies into three Conda prefixes. Each uses Python 3.12.16; retrieval includes Node.js 24.22.1, and publication conversion pins Pandoc 3.13. Conda package identities, pip inventories, editable project installation and official MCP commit are recorded separately. All three environments were rebuilt into fresh prefixes and passed dependency and version checks. A fixed non-editable VCS installation is identified by installation metadata; an editable installation additionally requires inspection of the checkout and tracked modifications.

% block-id: auto.07.platform
\hypertarget{auto.07.platform}{}
The observed desktop and execution host run Windows Server 2022 Datacenter, version 21H2, build 20348.4405. Codex, Chrome and Explorer were observed in session 1. Windows 10 remains the requested deployment target; the operating system of a possible remote client was not established by the host observation. This validation therefore applies to the observed Server 2022 installation. Windows 10 and Windows 11 require their own device evidence before platform support is claimed.

% block-id: auto.07.switch
\hypertarget{auto.07.switch}{}
Candidate MCP configurations are probed before replacing active configuration. After backup, the project and Codex configurations were switched and the MCP services restarted. Native tool responses confirmed the new Conda executable paths, Wolfram kernel 15.0.0 and twelve passing analytical tests. The ANSYS native client connected to a task-owned Student 2025.2.4 session using PyAEDT 1.8.0; disconnect and process cleanup were checked separately using recorded process creation identity. These checks establish connection and execution behavior for this device.

% block-id: auto.07.student
\hypertarget{auto.07.student}{}
An independent Python 3.13.18 environment reproduced missing Student-edition parameter propagation in the current PyAEDT-MCP upstream. The patch forwards the option through launch, connection and automatic server startup while preserving the default behavior. Nine added regressions failed on the baseline; the patched suite passed 141 tests. Actual Student launch, connection and startup calls also passed, with owned-process cleanup recorded. Pull request 155 was submitted to the official repository and its target branch and diff were read back. Submission is the observed contribution outcome; upstream review and merging are separate events. \cite{PYAEDTMCP}

% block-id: auto.07.research
\hypertarget{auto.07.research}{}
Four skills organize research discovery, paper acquisition, capability reuse and research handoff. Broad discovery prefers Tencent Yuanbao; source verification uses publisher pages, academic metadata, repository APIs and the authorized browser. The pinned MIT-licensed academic-search modules normalize records, retain metadata conflicts and acquire open-access files. Existing upstream identity, copyright and tests are preserved. Retrieval access, identity verification, software license and full-text redistribution rights are separate fields. \cite{ACADEMICSEARCH}

% block-id: auto.07.acquisition
\hypertarget{auto.07.acquisition}{}
The observed acquisition cases include a Chinese semi-analytical Litz-resistance paper and an English winding-loss paper. PDF parsing, title and DOI matching, and SHA256 establish the downloaded file identities. The Chinese publisher and institutional index expose different publication dates, which remain attributable observations after deduplication. The English accepted manuscript is available from White Rose under CC BY 4.0; the Chinese full text is retained locally with redistribution rights unresolved. The authorized institutional Chrome/WebVPN session subsequently downloaded a six-page PDF through the HTML reader. Its title and DOI passed identity checks; the institutional and publisher copies retain separate hashes. The history also preserves a blocked download, an alternate session requiring authentication, and human verification interventions. \cite{C01,ZHU2026}

% block-id: auto.07.handoff
\hypertarget{auto.07.handoff}{}
The handoff example saves a task description and parameters before running the existing round-wire CLI calibration. Twelve checks passed, including independent complex-current and Joule-power integrals. This validates a literature-to-task handoff to an existing model; reproducing the cited winding models requires additional geometry, local-field and convergence work. The small Maxwell replay passed port-current and power checks but retains its single-mesh limitation. A successful tool call and these limited physical checks do not establish a converged Litz performance ranking.

% block-id: auto.07.evaluation
\hypertarget{auto.07.evaluation}{}
Three paired trials compared live Crossref retrieval with reuse of a hash-verified same-day record for one DOI, using the same deterministic validator and permissions. Median times were 1.102 s and 0.344 ms; network calls were one and zero per trial, respectively. All six trials passed DOI, title, author and publication-date checks. The measurements exclude cache creation and human review. No model ran within these stages, and host tokens were unobservable. Broader efficiency claims require matched end-to-end tasks, model settings and tool permissions, with independent-device validation. LaTeX supplies the manuscript, with strict Pandoc checks on paragraph identities, prose and mathematics.
